% Generated draft: requires mathematical and source review.
\documentclass[11pt]{article}
\usepackage[margin=25mm]{geometry}
\usepackage{fontspec}
\setmainfont{Latin Modern Roman}
\usepackage{amsmath,amssymb,mathtools,unicode-math}
\setmathfont{Latin Modern Math}
\usepackage{enumitem,needspace,xurl,hyperref}
\setlength{\parindent}{0pt}
\setlength{\parskip}{5pt}
\setlength{\emergencystretch}{3em}
\newcommand{\N}{\mathbb{N}}
\newcommand{\Z}{\mathbb{Z}}
\newcommand{\Q}{\mathbb{Q}}
\newcommand{\R}{\mathbb{R}}
\newcommand{\C}{\mathbb{C}}
\newcommand{\F}{\mathbb{F}}
\newcommand{\A}{\mathbb{A}}
\newcommand{\PP}{\mathbb{P}}
\newcommand{\E}{\mathbb{E}}
\newcommand{\eps}{\varepsilon}
\newcommand{\dd}{\,\mathrm{d}}
\begin{document}
\section*{UnsolvedMath ID: 30000992}
\textbf{Source problem code:} catalogue entry
\subsection{Definitions and mathematical statement}
The source proposes a universality conjecture for grain growth. For generic initial conditions, let $\mathcal G_t$ denote the cell complex formed by grains evolving by the intended mean-curvature dynamics. After rescaling by the increasing average grain size, does this complex approach a universal random cell complex independent of the material?

A conditional distributional schema is useful to expose the missing data. Fix an intended dimension and class of materials, a law or spatial-sampling convention for initial conditions and observations, a length scale $L(t)$ derived from mean grain size, and a topology on the resulting space of cell complexes. Then the proposed conclusion would take the form
\[
\operatorname{Law}\bigl(L(t)^{-1}\mathcal G_t\bigr)
\Longrightarrow \nu_\infty\qquad(t\to\infty),
\]
for every material and generic initial condition in the specified classes, with the same limiting probability law $\nu_\infty$. Here the rescaling acts on the geometric realization of the complex. If average grain size means volume in dimension $d$, the associated length scale would conventionally be its $d$-th root; if it means a length, this conversion is unnecessary. The source does not choose either definition.

Neither randomness nor convergence is defined by the record. Deterministic evolution of a fixed complex does not by itself specify a law; sampling a typical grain and sampling a typical spatial point can give different statistics. The dimension, material parameters and their permitted variation, genericity class, mean-curvature and junction laws, continuation through singular events, boundary conditions, and finite-volume versus infinite-volume regime are also unspecified. In a bounded sample, the order of large-volume and long-time limits may be part of the required formulation. No one of these choices is selected by the displayed schema.

The supplied assessment checked 2026-08-21 expressly says that the universality conjecture is unproved and lacks a precise mathematical formulation; it also describes related later universality limits as conjectural. The open registry label and corrected-verified editorial-cleanup label are preserved as supplied metadata, not as evidence that the missing formulation has been completed.
\subsection{Short English statement}
Ask whether mean-curvature grain growth, rescaled by mean grain size, approaches one material-independent random cell-complex law for generic initial data. The source itself says that the probabilistic model, scale and convergence statement have not been made precise.
\subsection{Sources}
\begin{itemize}
\item Local UnsolvedMath registry, permanent ID 30000992, source code catalogue entry; supplied complete source record, including its editorial-cleanup note and explicit lack-of-precise-formulation assessment.
\item Computational Algebraic Topology (2009), Oberwolfach Reports, DOI 10.4171/owr/2008/29, pp. 1603--1654: \url{https://doi.org/10.4171/owr/2008/29}. Proposed-year metadata gives 2008.
\item Later coarsening universality study: \url{https://arxiv.org/abs/1505.07893}.
\end{itemize}
\end{document}
